% =====================================================================
% Fluid Dynamics and IBM Theory: Eulerian--Lagrangian coupling and regularized kernels
% =====================================================================
\documentclass[aspectratio=169,11pt]{beamer}

% Shared style, ported from the UF Animal Sciences 2026 talk (ib-tutorials.sty at the repo root).
\makeatletter\def\input@path{{./}{../}}\makeatother
\usepackage{ib-tutorials}

\usepackage[utf8]{inputenc}
\usepackage{graphicx}
\usepackage{amsmath,amssymb,bm,mathtools}
\usepackage{booktabs}
\usepackage{xcolor}
\usepackage{tcolorbox}
\tcbuselibrary{skins}
\usepackage{hyperref}

\graphicspath{{../pictures/1-2/}{pictures/1-2/}}

\newcommand{\xvec}{\bm{x}}
% Module 1 uses the Peskin convention throughout:
%   \Xvec   = the material (Lagrangian) label  q
%   \chivec = the CURRENT POSITION of that material point, X(q,t)
% Module 4 (IBFE) instead writes x = \chi(X,t) with X the reference point;
% that notation appears in this deck only on the explicit bridge slide.
\newcommand{\Xvec}{\bm{q}}
\newcommand{\chivec}{\bm{X}}
\newcommand{\uvec}{\bm{u}}
\newcommand{\Uvec}{\bm{U}}
\newcommand{\fvec}{\bm{f}}
\newcommand{\Fvec}{\bm{F}}
\newcommand{\Pten}{\bm{P}_{e}}
\newcommand{\Nvec}{\bm{N}}
\newcommand{\testvec}{\bm{\varphi}}

\title[IB coupling and kernels]{Integral transforms and regularized delta functions}
\subtitle{Force transfer, material measures, and discrete power consistency}
\author[Miller]{Laura A. Miller}
\institute[UArizona]{Department of Mathematics, University of Arizona}
\date{Fluid Dynamics and IBM Theory \textbullet\ Module 1}

\begin{document}

% ----- 1. Title -----
% The title slide uses the shared \maketutorialtitlepanel macro from
% ib-tutorials.sty: same layout as \maketutorialtitle, with the right column
% holding this deck's formula panel instead of an image.
\begin{frame}[plain]
\maketutorialtitlepanel{%
    \begin{punchbox}
    \centering
    \[
      \begin{aligned}
      \Fvec(\Xvec,t)&\xrightarrow{\ \text{spread}\ }\fvec(\xvec,t),\\[3pt]
      \uvec(\xvec,t)&\xrightarrow{\ \text{interpolate}\ }\Uvec(\Xvec,t).
      \end{aligned}
    \]
    \small One kernel links two coordinate descriptions.
    \end{punchbox}}
\end{frame}

% ----- 2. Roadmap -----
\begin{frame}{Roadmap}
\begin{columns}[T,onlytextwidth]
  \begin{column}{0.32\textwidth}
  \begin{block}{1. Continuous coupling}
    \small Reference and spatial coordinates, force spreading, velocity interpolation
  \end{block}
  \end{column}
  \begin{column}{0.32\textwidth}
  \begin{block}{2. Material measures}
    \small The boundary traction term, unified versus partitioned forces, and the Jacobian question
  \end{block}
  \end{column}
  \begin{column}{0.32\textwidth}
  \begin{block}{3. Discrete kernels}
    \small Moment conditions, approximate translation invariance, and the power identity
  \end{block}
  \end{column}
\end{columns}

\vfill
\begin{punchbox}
\centering
The central modeling choice is the \hi{measure with respect to which a force density is defined}.
\end{punchbox}
\end{frame}

\section{Continuous coupling}

% ----- 3. Coordinates -----
\begin{frame}{Two coordinate descriptions of one system}
\begin{columns}[T,onlytextwidth]
  \begin{column}{0.48\textwidth}
    \begin{block}{Eulerian description}
      \[
        \xvec\in\Omega\subset\mathbb{R}^{d},
        \qquad
        \uvec=\uvec(\xvec,t),\quad p=p(\xvec,t).
      \]
      The fluid variables live on a fixed spatial domain.
    \end{block}
  \end{column}
  \begin{column}{0.48\textwidth}
    \begin{block}{Lagrangian description}
      \[
        \Xvec\in\mathcal U\subset\mathbb{R}^{k},
        \qquad
        \xvec=\chivec(\Xvec,t),\quad k\le d.
      \]
      The label \(\Xvec\) identifies a material point.
    \end{block}
  \end{column}
\end{columns}

\medskip
\begin{punchbox}
\(\mathcal U\) is a \hi{fixed material-coordinate domain}. It may be chosen as a reference configuration, but it need not equal the initial spatial configuration: \(\chivec(\Xvec,0)=\Xvec\) is not required (for a fiber or surface, \(k<d\), it is not even dimensionally possible).
\end{punchbox}

\smallskip
\soft{Some texts write \(s\) or \(\bm\xi\) for the same material coordinate. Module 4 (IBFE) switches to continuum-mechanics notation, \(\bm\chi(\bm X,t)\) with \(\bm X\) the reference point; the next slide compares the two.}
\end{frame}

% ----- Notation bridge -----
\begin{frame}{Two equivalent Lagrangian notations}
\begin{columns}[T,onlytextwidth]
  \begin{column}{0.48\textwidth}
  \begin{block}{Classical IB notation (used in this module)}
  \[
    \xvec=\bm{X}(\bm q,t),
    \qquad \bm q\in\mathcal U .
  \]
  The symbol \(\bm q\) labels a point on a curve, surface, or material mesh; \(\bm X\) is where that point currently \emph{is}.
  \end{block}
  \end{column}
  \begin{column}{0.48\textwidth}
  \begin{block}{Continuum-mechanics notation (used from Module 4 on)}
  \[
    \xvec=\bm{\chi}(\bm{X},t),
    \qquad \bm{X}\in\mathcal B_0 .
  \]
  Here \(\bm X\) is the \emph{reference} point in the reference configuration \(\mathcal B_0\). This notation makes reference derivatives and stress measures explicit.
  \end{block}
  \end{column}
\end{columns}

\medskip
\begin{punchbox}
These describe the same material-to-spatial map, and the roles of the symbol \(\bm X\) are swapped between them. \warn{In this module \(\bm q\) is the material label and \(\bm X(\bm q,t)\) its current position}; Module 4 (IBFE) switches to \(\bm\chi(\bm X,t)\) with \(\bm X\) the reference point.
\end{punchbox}
\end{frame}

% ----- Material domain -----
\begin{frame}{What does the material domain represent?}
\begin{itemize}\itemsep7pt
  \item \(\mathcal U\) is fixed in time and indexes the structural degrees of freedom.
  \item Integration over \(\mathcal U\) sums \hi{material points}, not points of the Eulerian grid.
  \item Its dimension \(k\) may differ from the ambient dimension \(d\): fibers have \(k=1\), surfaces have \(k=2\), and volumetric solids have \(k=d\).
  \item Geometry enters through the current placement \(\chivec(\mathcal U,t)\subset\Omega\).
\end{itemize}

\[
  \boxed{\mathcal U\ \neq\ \chivec(\mathcal U,t)\ \subset\ \Omega .}
\]

\begin{punchbox}
Calling \(\mathcal U\) a ``reference domain'' does not require \(\chivec(\Xvec,0)=\Xvec\). A convenient material parameterization need not coincide with the initial spatial geometry.
\end{punchbox}
\end{frame}

% ----- 4. Visual coupling -----
\begin{frame}{Coupling goes in both directions}
\begin{columns}[T,onlytextwidth]
  \begin{column}{0.48\textwidth}
    \centering
    \includegraphics[height=0.50\textheight]{force-spread2.png}

    \hi{Spread force:} structure \(\longrightarrow\) fluid grid

    {\scriptsize Figure from deck 1-2.}\\
    {\tiny\color{DarkGray}Schematic: instructor's course notes.}
  \end{column}
  \begin{column}{0.48\textwidth}
    \centering
    \includegraphics[height=0.50\textheight]{vel-spread2.png}

    \hi{Interpolate velocity:} fluid grid \(\longrightarrow\) structure

    {\scriptsize Figure from deck 1-2.}\\
    {\tiny\color{DarkGray}Schematic: instructor's course notes.}
  \end{column}
\end{columns}

\medskip
\begin{punchbox}
\centering
The continuum Dirac delta locates the interaction exactly; a discrete kernel shares it among nearby grid points.
\end{punchbox}
\end{frame}

% ----- 5. Continuous transforms -----
\begin{frame}{The two continuous interaction equations}
\small
Here \(\Xvec\) is the fixed material label, \(\chivec(\Xvec,t)\) its current position, and \(\Fvec(\Xvec,t)\) force per unit material measure \(d\Xvec\) \soft{(Peskin convention, as in decks 1-1, 1-2, 1-4)}.
\normalsize

\begin{block}{Force spreading: Lagrangian \(\rightarrow\) Eulerian}
\[
  \boxed{
  \fvec(\xvec,t)
  =
  \int_{\mathcal U}
  \Fvec(\Xvec,t)\,
  \delta\!\left(\xvec-\chivec(\Xvec,t)\right)\,d\Xvec }
\]
\end{block}

\begin{block}{Velocity interpolation: Eulerian \(\rightarrow\) Lagrangian}
\[
  \boxed{
  \Uvec(\Xvec,t)
  =
  \int_{\Omega}
  \uvec(\xvec,t)\,
  \delta\!\left(\xvec-\chivec(\Xvec,t)\right)\,d\xvec }
\]
\end{block}

For a no-slip material structure, \(\partial_t\chivec=\Uvec\); with the true Dirac delta this is equivalent to \(\partial_t\chivec(\Xvec,t)=\uvec(\chivec(\Xvec,t),t)\).
\end{frame}

% ----- Force spreading in detail -----
\begin{frame}{Force spreading is a material-to-spatial map}
Let \(\Fvec(\Xvec,t)\,d\Xvec\) be the force carried by a material element. Its Eulerian image is
\[
  \boxed{
  \fvec(\xvec,t)
  =\int_{\mathcal U}\Fvec(\Xvec,t)
  \delta\!\left(\xvec-\chivec(\Xvec,t)\right)\,d\Xvec .}
\]

\begin{enumerate}\itemsep7pt
  \item Compute the structural force in material coordinates.
  \item Use \(\chivec(\Xvec,t)\) to locate the material element in physical space.
  \item Deposit its force at that location through the Dirac distribution.
\end{enumerate}

\begin{punchbox}
The integral defines a force \emph{measure} on \(\Omega\). It remains meaningful when the structure has lower dimension than the fluid domain.
\end{punchbox}
\end{frame}

% ----- Velocity interpolation in detail -----
\begin{frame}{Velocity interpolation closes the kinematic coupling}
The reverse transform samples the Eulerian velocity at the current material position:
\[
  \Uvec(\Xvec,t)
  =\int_{\Omega}\uvec(\xvec,t)
  \delta\!\left(\xvec-\chivec(\Xvec,t)\right)\,d\xvec
  =\uvec\!\left(\chivec(\Xvec,t),t\right).
\]

For a no-slip immersed structure,
\[
  \boxed{\partial_t\chivec(\Xvec,t)=\Uvec(\Xvec,t).}
\]

\begin{columns}[T,onlytextwidth]
  \begin{column}{0.48\textwidth}
  \begin{block}{Dynamic coupling}
  \(\Fvec\mapsto\fvec\): the structure loads the fluid.
  \end{block}
  \end{column}
  \begin{column}{0.48\textwidth}
  \begin{block}{Kinematic coupling}
  \(\uvec\mapsto\partial_t\chivec\): the fluid moves the structure.
  \end{block}
  \end{column}
\end{columns}
\end{frame}

% ----- 6. Weak meaning -----
\begin{frame}{What the force-spreading formula means}
The Eulerian force may be a distribution rather than an ordinary function. Its precise meaning is: for every smooth, compactly supported test field \(\testvec\),
\[
  \boxed{
  \int_{\Omega}
  \fvec(\xvec,t)\cdot\testvec(\xvec)\,d\xvec
  =
  \int_{\mathcal U}
  \Fvec(\Xvec,t)\cdot
  \testvec\!\left(\chivec(\Xvec,t)\right)\,d\Xvec .}
\]

\begin{columns}[T,onlytextwidth]
  \begin{column}{0.48\textwidth}
  \begin{block}{Constant test field}
    Total Eulerian force equals total Lagrangian force.
  \end{block}
  \end{column}
  \begin{column}{0.48\textwidth}
  \begin{block}{Rigid-rotation test field}
    The first spatial moment gives the corresponding torque identity.
  \end{block}
  \end{column}
\end{columns}

\begin{punchbox}
The delta formula is a \hi{distributional push-forward of force measure}; it is not a change of integration variables.
\end{punchbox}
\end{frame}

% ----- Region test, steps 1--2 -----
\begin{frame}{Test force spreading on a spatial region: steps 1--2}
Choose a region \(\Gamma\subset\Omega\). Formally, its total Eulerian force is
\[
  \int_{\Gamma}\fvec(\xvec,t)\,d\xvec .
\]
Substitute the spreading relation:
\[
  \int_{\Gamma}\fvec(\xvec,t)\,d\xvec
  =
  \int_{\Gamma}\int_{\mathcal U}
  \Fvec(\Xvec,t)\,
  \delta\!\left(\xvec-\chivec(\Xvec,t)\right)
  \,d\Xvec\,d\xvec .
\]

\begin{punchbox}
This notation asks a geometric question: which material points are currently located inside \(\Gamma\)?
\end{punchbox}
\end{frame}

% ----- Region test, steps 3--4 -----
\begin{frame}{Test force spreading on a spatial region: steps 3--4}
After moving the material integral outside, the inner delta action is
\[
  \int_{\Gamma}\delta\!\left(\xvec-\chivec(\Xvec,t)\right)\,d\xvec
  =
  \begin{cases}
  1, & \chivec(\Xvec,t)\in\Gamma,\\
  0, & \chivec(\Xvec,t)\notin\Gamma,
  \end{cases}
\]
away from the usual boundary ambiguity.

\medskip
Thus only material points whose \emph{current} positions lie in \(\Gamma\) contribute.

\begin{punchbox}
The material integration set is the preimage
\(
\chivec_t^{-1}(\Gamma)
=\{\Xvec\in\mathcal U:\chivec(\Xvec,t)\in\Gamma\}.
\)
\end{punchbox}
\end{frame}

% ----- Region identity -----
\begin{frame}{The region identity is the physical punchline}
Combining the preceding steps gives the formal identity
\[
  \boxed{
  \int_{\Gamma}\fvec(\xvec,t)\,d\xvec
  =
  \int_{\chivec_t^{-1}(\Gamma)}\Fvec(\Xvec,t)\,d\Xvec .}
\]

\begin{block}{Interpretation}
The total force assigned to a fluid region equals the total force carried by material points occupying that region.
\end{block}

\begin{punchbox}
This statement explains force spreading without imagining a point force at every grid node. The continuum object is a pushed-forward force measure.
\end{punchbox}
\end{frame}

% ----- Fubini caveat -----
\begin{frame}{Can we really swap the integrals?}
\begin{columns}[T,onlytextwidth]
  \begin{column}{0.48\textwidth}
  \begin{block}{True Dirac delta}
  The delta is a distribution, not an \(L^1\) function. Classical Fubini--Tonelli is therefore not the right justification for the formal double integral.

  Define \(\fvec\) by its action on smooth test fields and approximate the indicator of \(\Gamma\).
  \end{block}
  \end{column}
  \begin{column}{0.48\textwidth}
  \begin{block}{Regularized kernel \(\delta_h\)}
  If \(\delta_h\) and \(\Fvec\) are integrable, ordinary Fubini applies. A marker within \(O(h)\) of \(\partial\Gamma\) then contributes only partially to the region.
  \end{block}
  \end{column}
\end{columns}
\vspace{4pt}
\begin{punchbox}
The rigorous continuum starting point is the smooth-test-field identity on the slide ``What the force-spreading formula means'' (four slides back), not the region identity on the previous slide.
\end{punchbox}
\end{frame}

\section{Material measures and elastic stress}

% ----- 7. Jacobian -----
\begin{frame}{Where did the Jacobian go?}
\begin{columns}[T,onlytextwidth]
  \begin{column}{0.48\textwidth}
  \begin{block}{Reference-measure density}
  If \(\Fvec_{\mathrm{ref}}\) is force per \(d\Xvec\), then
  \[
    \fvec(\xvec)
    =\int_{\mathcal U}\Fvec_{\mathrm{ref}}(\Xvec)
    \delta(\xvec-\chivec(\Xvec))\,d\Xvec .
  \]
  No variable substitution is being made.
  \end{block}
  \end{column}
  \begin{column}{0.48\textwidth}
  \begin{block}{Current-measure density}
  For a volumetric map, \(d\xvec=J\,d\Xvec\), so the same force satisfies
  \[
    \Fvec_{\mathrm{ref}}=J\Fvec_{\mathrm{cur}},
    \qquad
    J=\det\!\left(\frac{\partial\chivec}{\partial\Xvec}\right).
  \]
  Hence
  \[
    \fvec=\int_{\mathcal U}\Fvec_{\mathrm{cur}}
    \delta(\xvec-\chivec)J\,d\Xvec .
  \]
  \end{block}
  \end{column}
\end{columns}

\begin{punchbox}
\centering
No Jacobian is ``missing'': it appears only when the chosen force density is expressed per unit \hi{current} measure.
\end{punchbox}
\end{frame}

% ----- Metric factors -----
\begin{frame}{When does a Jacobian or metric factor appear?}
Suppose the force density is specified per unit \emph{current} geometric measure. Converting it to the fixed material measure introduces the appropriate stretch:

\begin{center}
\begin{tabular}{@{}lll@{}}
\toprule
Structure & Measure transformation & Reference/current density relation \\
\midrule
Volume, \(k=d\) & \(d\xvec=J\,d\Xvec\) & \(\Fvec_{\rm ref}=J\Fvec_{\rm cur}\) \\
Surface in 3D & \(da=J_s\,dA\) & \(\Fvec_{\rm ref}=J_s\Fvec_{\rm cur}\) \\
Fiber & \(ds=\lambda\,dS\) & \(\Fvec_{\rm ref}=\lambda\Fvec_{\rm cur}\) \\
\bottomrule
\end{tabular}
\end{center}

\begin{punchbox}
The factor belongs to a \hi{density conversion}. It is not generated merely because \(\chivec\) appears inside a delta distribution.
\end{punchbox}
\end{frame}

% ----- Same physical force -----
\begin{frame}{Reference and current densities represent the same force}
For a volumetric placement, the force carried by one material element satisfies
\[
  \Fvec_{\rm ref}(\Xvec,t)\,d\Xvec
  =\Fvec_{\rm cur}(\Xvec,t)\,d\xvec,
  \qquad d\xvec=J\,d\Xvec .
\]
Therefore
\[
  \boxed{\Fvec_{\rm ref}=J\Fvec_{\rm cur}}
\]
and the same Eulerian force can be written either way:
\[
  \fvec
  =\int_{\mathcal U}\Fvec_{\rm ref}\,\delta(\xvec-\chivec)\,d\Xvec
  =\int_{\mathcal U}\Fvec_{\rm cur}\,\delta(\xvec-\chivec)J\,d\Xvec .
\]

\begin{punchbox}
Always report both the density and its underlying measure; otherwise two equivalent formulas can appear contradictory.
\end{punchbox}
\end{frame}

% ----- Elastic stress in material coordinates -----
\begin{frame}{Elasticity is evaluated in material coordinates}
For a volumetric hyperelastic solid, define
\[
  \bm F_{\rm def}(\Xvec,t)=\frac{\partial\chivec}{\partial\Xvec},
  \qquad
  J=\det\bm F_{\rm def}>0 .
\]
If \(W_e(\bm F_{\rm def})\) is energy per unit reference volume, the first Piola--Kirchhoff stress is
\[
  \boxed{
  \Pten(\Xvec,t)
  =\frac{\partial W_e}{\partial\bm F_{\rm def}} .}
\]

\begin{itemize}\itemsep6pt
  \item \(\Pten\) maps a reference normal to force per unit reference area.
  \item Its material divergence has units of force per unit reference volume.
  \item \(\Pten\) is generally nonsymmetric, even when the Cauchy stress is symmetric.
\end{itemize}
\end{frame}

% ----- 8. Strong elastic force -----
\begin{frame}{From elastic stress to Eulerian force}
The material stress produces the interior force density
\[
  \bm G=\operatorname{Div}_{\!\Xvec}\Pten
  \qquad\text{in }\mathcal U,
\]
and the solid boundary transmits the reaction \(-\Pten\Nvec\) to the surrounding fluid.

\begin{block}{Strong-form elastic force (Griffith--Luo)}
\[
\boxed{
\begin{aligned}
\fvec_e(\xvec,t)
={}&\int_{\mathcal U}
\operatorname{Div}_{\!\Xvec}\Pten(\Xvec,t)\,
\delta(\xvec-\chivec(\Xvec,t))\,d\Xvec \\
&-\int_{\partial\mathcal U}
\Pten(\Xvec,t)\Nvec(\Xvec)\,
\delta(\xvec-\chivec(\Xvec,t))\,dA(\Xvec).
\end{aligned}}
\]
\end{block}

\small The first integral spreads the internal stress divergence; the second spreads the boundary transmission force. The second term may be omitted only when its traction vanishes or when boundary effects are represented another way.
\end{frame}

% ----- 9. Internal versus transmission -----
\begin{frame}{Two elastic-force contributions}
\begin{columns}[T,onlytextwidth]
  \begin{column}{0.48\textwidth}
  \begin{block}{Internal force density}
    \[
      \bm{G}=\operatorname{Div}_{\!\Xvec}\Pten
      \qquad\text{in }\mathcal U.
    \]
    \begin{itemize}
      \item per unit reference volume
      \item supported throughout the solid
      \item regular when \(\Pten\) is smooth
    \end{itemize}
  \end{block}
  \end{column}
  \begin{column}{0.48\textwidth}
  \begin{block}{Transmission force density}
    \[
      \bm{T}=-\Pten\Nvec
      \qquad\text{on }\partial\mathcal U.
    \]
    \begin{itemize}
      \item per unit reference area
      \item supported on the solid boundary
      \item produces a singular Eulerian force layer
    \end{itemize}
  \end{block}
  \end{column}
\end{columns}

\begin{punchbox}
For a lower-dimensional interface, the whole force is singular in the ambient fluid. For a volumetric solid, only the transmitted boundary force is confined to a lower-dimensional set.
\end{punchbox}
\end{frame}

% ----- Unified weak formulation -----
\begin{frame}{Unified weak force: one material density}
Define \(\Fvec\) through virtual work: for every structural test field \(\bm V\),
\[
  \boxed{
  \int_{\mathcal U}\Fvec\cdot\bm V\,d\Xvec
  =-\int_{\mathcal U}\Pten:\nabla_{\!\Xvec}\bm V\,d\Xvec .}
\]
The test field need not vanish on \(\partial\mathcal U\). Integration by parts gives
\[
  \int_{\mathcal U}\Fvec\cdot\bm V\,d\Xvec
  =\int_{\mathcal U}(\operatorname{Div}\Pten)\cdot\bm V\,d\Xvec
  -\int_{\partial\mathcal U}(\Pten\Nvec)\cdot\bm V\,dA .
\]

\begin{punchbox}
The projected volume field \(\Fvec\) packages both the interior divergence and the boundary transmission force. No traction-free assumption is being made.
\end{punchbox}
\end{frame}

% ----- Partitioned weak formulation -----
\begin{frame}{Partitioned weak force: volume and boundary fields}
Define the interior density \(\bm G\) by
\[
  \int_{\mathcal U}\bm G\cdot\bm V\,d\Xvec
  =-\int_{\mathcal U}\Pten:\nabla_{\!\Xvec}\bm V\,d\Xvec
  +\int_{\partial\mathcal U}(\Pten\Nvec)\cdot\bm V\,dA,
\]
so that \(\bm G=\operatorname{Div}\Pten\) weakly, and set
\[
  \bm T=-\Pten\Nvec\qquad\text{on }\partial\mathcal U .
\]
Then spread the two contributions separately:
\[
  \fvec
  =\int_{\mathcal U}\bm G\,\delta(\xvec-\chivec)\,d\Xvec
  +\int_{\partial\mathcal U}\bm T\,\delta(\xvec-\chivec)\,dA .
\]
\end{frame}

% ----- 10. Unified versus partitioned -----
\begin{frame}{Unified and partitioned weak formulations (side-by-side recap)}
\small
\soft{The previous two slides, placed side by side for comparison; the formulas are the same.}\par\smallskip
\begin{columns}[T,onlytextwidth]
  \begin{column}{0.48\textwidth}
  \begin{block}{Unified force}
  Define one density \(\Fvec\) weakly by
  \[
    \int_{\mathcal U}\Fvec\cdot\bm{V}\,d\Xvec
    =-\int_{\mathcal U}\Pten:\nabla_{\!\Xvec}\bm{V}\,d\Xvec,
  \]
  for test fields \(\bm{V}\) that need not vanish on \(\partial\mathcal U\), and spread
  \[
    \fvec=\int_{\mathcal U}\Fvec\,\delta(\xvec-\chivec)\,d\Xvec.
  \]
  Integration by parts shows that \(\Fvec\) packages both contributions.
  \end{block}
  \end{column}
  \begin{column}{0.48\textwidth}
  \begin{block}{Partitioned force}
  Keep \(\bm{G}\) and \(\bm{T}\) separate:
  \[
    \begin{aligned}
    \int_{\mathcal U}\bm{G}\cdot\bm{V}\,d\Xvec
    ={}&-\int_{\mathcal U}\Pten:\nabla_{\!\Xvec}\bm{V}\,d\Xvec\\
    &+\int_{\partial\mathcal U}\Pten\Nvec\cdot\bm{V}\,dA,
    \end{aligned}
  \]
  \[
    \fvec=\underbrace{\int_{\mathcal U}\bm{G}\,\delta\,d\Xvec}_{\bm{g}}
    +\underbrace{\int_{\partial\mathcal U}\bm{T}\,\delta\,dA}_{\bm{t}},
    \qquad \bm{T}=-\Pten\Nvec.
  \]
  \end{block}
  \end{column}
\end{columns}
\end{frame}

% ----- Units example -----
\begin{frame}{Units check: an immersed curve in a 2-D fluid}
Let the material label \(q\) be arclength along a one-dimensional structure, measured in its reference configuration so that the label is fixed in time, so that \(\bm X(q)\) is the current position of that material point, and let the two-dimensional fluid force density be
\[
  \fvec(\xvec)=\int\Fvec(q)\,\delta\!\left(\xvec-\bm X(q)\right)\,dq .
\]
Then
\[
  [\Fvec]=\frac{\mathrm N}{\mathrm m},
  \qquad
  [\delta]=\frac{1}{\mathrm m^2},
  \qquad
  [dq]=\mathrm m,
\]
and therefore
\[
  [\fvec]
  =\frac{\mathrm N}{\mathrm m}\frac{1}{\mathrm m^2}\mathrm m
  =\boxed{\frac{\mathrm N}{\mathrm m^2}}.
\]

\begin{punchbox}
In a 2-D model this is force per unit area. If the 2-D model is read as a slice of an extruded 3-D problem, every force is additionally per unit out-of-plane depth: \([\Fvec]=\mathrm N/\mathrm m^2\) and \([\fvec]=\mathrm N/\mathrm m^3\), the familiar 3-D body-force density.
\end{punchbox}
\end{frame}

% ----- 11. Units -----
\begin{frame}{Units follow the chosen material measure}
Assume first that the material coordinates have units of length. For a \(k\)-dimensional structure in a \(d\)-dimensional fluid,
\[
  [\Fvec_{\mathrm{ref}}]=\frac{\mathrm{N}}{\mathrm{m}^{k}},
  \qquad
  [\delta]=\frac{1}{\mathrm{m}^{d}},
  \qquad
  [d\Xvec]=\mathrm{m}^{k},
\]
so \([\fvec]=\mathrm{N}/\mathrm{m}^{d}\).

\vspace{2pt}
{\centering
\begin{tabular}{@{}llll@{}}
\toprule
Structure & \(k\) & Reference measure & Force-density units \\
\midrule
Fiber & 1 & length & \(\mathrm{N}/\mathrm{m}\) \\
Surface & 2 & area & \(\mathrm{N}/\mathrm{m}^{2}\) \\
Volumetric solid & \(d\) & \(d\)-volume & \(\mathrm{N}/\mathrm{m}^{d}\) \\
\bottomrule
\end{tabular}\par}
\vspace{4pt}
\begin{punchbox}
If the material labels are dimensionless, \(\Fvec\) instead has units of force per unit \hi{parameter measure}; any geometric scale factors are already part of that density or of the quadrature weight.
\end{punchbox}
\end{frame}

% ----- Distributional force layer -----
\begin{frame}{The continuum force can be a singular layer}
For a lower-dimensional immersed structure, \(\fvec\) is a distribution supported on \(\chivec(\mathcal U,t)\), not an ordinary pointwise field.

\begin{itemize}\itemsep7pt
  \item A curve in 2-D or a surface in 3-D creates a delta-supported force layer.
  \item Such a layer can produce jumps in pressure or viscous traction across the interface.
  \item Under classical no-slip coupling, the fluid velocity itself remains continuous across a material interface, although its derivatives may jump.
\end{itemize}

\begin{punchbox}
``Point force'' is misleading here: each material element carries finite force \(\Fvec\,d\Xvec\), and the collection forms a singular measure along the immersed geometry.
\end{punchbox}
\end{frame}

\section{Regularized delta functions}

% ----- 12. Regularization -----
\begin{frame}{From a Dirac delta to a grid kernel}
Conventional IB discretizations replace the singular delta distribution by a compactly supported kernel tied to the Eulerian spacing \(h\):
\[
  \boxed{
  \delta_h(\xvec)
  =\frac{1}{h^{d}}
  \prod_{\alpha=1}^{d}
  \phi\!\left(\frac{x_{\alpha}}{h}\right).}
\]

\begin{columns}[T,onlytextwidth]
  \begin{column}{0.48\textwidth}
  \begin{block}{Why regularize?}
    \begin{itemize}
      \item Couple nonmatching grids locally
      \item Avoid placing a singular distribution on grid nodes
      \item Obtain sparse interaction operators
    \end{itemize}
  \end{block}
  \end{column}
  \begin{column}{0.48\textwidth}
  \begin{block}{What tensor products imply}
    \begin{itemize}
      \item Efficient, separable evaluation
      \item Multidimensional moments inherited from \(\phi\)
      \item Grid alignment: rotational invariance is only approximate
    \end{itemize}
  \end{block}
  \end{column}
\end{columns}

\soft{Sharp-interface and fully variational immersed methods need not use a regularized delta kernel.}
\end{frame}

% ----- 13. Discrete transforms -----
\begin{frame}{Discrete spreading and interpolation}
Let \(\xvec_{\bm{i}}\) be Eulerian grid points, \(\chivec_m=\chivec(\Xvec_m)\), and \(w_m\) Lagrangian quadrature weights.

\begin{block}{Matched interaction operators}
\begin{align*}
  \fvec_{\bm{i}}
  &=\sum_m \Fvec_m\,
  \delta_h(\xvec_{\bm{i}}-\chivec_m)\,w_m,
  &&\text{(spread)}\\[2mm]
  \Uvec_m
  &=\sum_{\bm{i}} \uvec_{\bm{i}}\,
  \delta_h(\xvec_{\bm{i}}-\chivec_m)\,h^d,
  &&\text{(interpolate)}
\end{align*}
\end{block}

\begin{punchbox}
The same kernel is necessary but not, by itself, the whole consistency statement. The Eulerian cell-volume factors and Lagrangian quadrature weights determine the two discrete inner products.
\end{punchbox}
\vspace{3pt}
\small A finite-element IB method may evaluate \(\chivec\) and \(\Fvec\) at element quadrature points rather than only at structural nodes.
\end{frame}

% ----- Kernel design criteria -----
\begin{frame}{Why does the regularized kernel need special properties?}
A generic smooth bump can spread a marker force, but it may not respect the discrete mechanics.

\begin{enumerate}\itemsep6pt
  \item \hi{Compact support} keeps coupling local and sparse.
  \item \hi{Discrete moments} reproduce constants and linear fields.
  \item \hi{Reflection symmetry} avoids a built-in left--right bias.
  \item A \hi{constant sum of squares} reduces sensitivity to subgrid translation.
  \item \hi{Smoothness} controls how interaction values vary as markers cross the grid.
\end{enumerate}

\begin{punchbox}
These are distinct algebraic conditions. Conservation and power consistency must be checked from the paired discrete operators, not inferred from a kernel's appearance.
\end{punchbox}
\end{frame}

% ----- Compact support -----
\begin{frame}{Property 1: compact support gives locality}
For a four-point one-dimensional kernel,
\[
  \phi(r)=0\qquad\text{when }|r|\ge2 .
\]
In \(d\) dimensions,
\[
  \delta_h(\xvec)
  =h^{-d}\prod_{\alpha=1}^{d}\phi(x_\alpha/h),
\]
so a Lagrangian sample interacts with only a small tensor-product stencil.

\begin{columns}[T,onlytextwidth]
  \begin{column}{0.48\textwidth}
  \begin{block}{Benefit}
  Sparse coupling, predictable cost, and communication confined to nearby grid cells.
  \end{block}
  \end{column}
  \begin{column}{0.48\textwidth}
  \begin{block}{Tradeoff}
  Wider support can buy smoother or higher-moment kernels, but increases work and blurs the interface over more cells.
  \end{block}
  \end{column}
\end{columns}
\end{frame}

% ----- Zeroth moment -----
\begin{frame}{Property 2: the zeroth moment preserves total force}
Require, for every real offset \(r\),
\[
  \boxed{\sum_{j\in\mathbb Z}\phi(r-j)=1 .}
\]
Then spreading a single material force \(\Fvec_m w_m\) gives
\[
  \sum_{\bm i}\fvec_{\bm i}h^d
  =\Fvec_m w_m
  \prod_{\alpha=1}^{d}\left(\sum_{i_\alpha}\phi(r_\alpha-i_\alpha)\right)
  =\Fvec_m w_m .
\]

\begin{punchbox}
The same identity makes interpolation reproduce constant Eulerian velocity fields exactly.
\end{punchbox}
\end{frame}

% ----- First moment -----
\begin{frame}{Property 3: the first moment centers the interaction}
Require
\[
  \boxed{\sum_{j\in\mathbb Z}(r-j)\phi(r-j)=0 .}
\]
Equivalently, a marker at \(X=rh\) satisfies
\[
  \sum_j (X-jh)\,\delta_h(X-jh)\,h=0 .
\]

\begin{itemize}\itemsep6pt
  \item The discrete force is centered at the marker rather than a nearby grid node.
  \item Together with the zeroth moment, spreading preserves the corresponding first spatial moment and hence total torque.
  \item Interpolation reproduces linear fields exactly.
\end{itemize}
\end{frame}

% ----- Symmetry and translation -----
\begin{frame}{Property 4: separate symmetry from translation invariance}
Even symmetry,
\[
  \phi(r)=\phi(-r),
\]
gives reflection symmetry of the one-dimensional stencil.

\medskip
Peskin's constant-sum-of-squares condition,
\[
  \sum_{j\in\mathbb Z}\phi(r-j)^2=C
  \qquad\text{independent of }r,
\]
reduces variation in a marker's self-coupling as it moves through one grid cell.

\begin{punchbox}
Neither condition makes a tensor-product kernel exactly rotationally or translationally invariant. They control specific grid-position effects.
\end{punchbox}
\end{frame}

% ----- 14. Moments -----
\begin{frame}{The first two discrete moment conditions}
\small
For every real offset \(r\), impose
\begin{align*}
  \sum_{j\in\mathbb Z}\phi(r-j)&=1,
  &\text{zeroth moment},\\
  \sum_{j\in\mathbb Z}(r-j)\phi(r-j)&=0,
  &\text{first moment}.
\end{align*}

\begin{columns}[T,onlytextwidth]
  \begin{column}{0.48\textwidth}
  \begin{block}{Zeroth moment}
    \begin{itemize}
      \item interpolates constants exactly
      \item preserves total force in spreading
    \end{itemize}
  \end{block}
  \end{column}
  \begin{column}{0.48\textwidth}
  \begin{block}{First moment}
    \begin{itemize}
      \item centers the discrete force
      \item gives total-torque agreement
      \item permits exact interpolation of linear fields
    \end{itemize}
  \end{block}
  \end{column}
\end{columns}

\begin{punchbox}
Together these identities give \(O(h^2)\) interpolation error for sufficiently smooth fields; they do not imply high-order accuracy at a singular fluid--structure interface.
\end{punchbox}
\end{frame}

% ----- 15. Other kernel properties -----
\begin{frame}{Keep the other kernel properties distinct (recap)}
\small
\soft{Recap of Properties 1 and 4, collected with smoothness so the four non-moment properties can be compared in one place.}
\begin{description}
  \item[Compact support] Localizes the \emph{coupling stencil}. The support radius is a kernel choice (for example, two mesh widths for a four-point kernel).
  \item[Even symmetry] \(\phi(r)=\phi(-r)\) gives reflection symmetry. It does not make a tensor-product kernel exactly rotationally invariant.
  \item[Sum of squares] Peskin's construction can impose
  \[
    \sum_{j\in\mathbb Z}\phi(r-j)^2=C,
    \qquad C\text{ independent of }r.
  \]
  This reduces variation in self-coupling as a marker moves relative to the grid and is an \emph{approximate translation-invariance} condition.
  \item[Smoothness] A smoother \(\phi\) produces smoother interpolated quantities, usually at the cost of a wider stencil or additional constraints.
\end{description}

\begin{punchbox}
\warn{Do not confuse the sum-of-squares condition with energy consistency.} Energy consistency comes from discrete adjointness.
\end{punchbox}
\end{frame}

% ----- 16. Power identity -----
\begin{frame}{Why the two operators should be adjoints}
Define the discrete inner products
\[
  (\uvec,\fvec)_h
  =\sum_{\bm{i}}\uvec_{\bm{i}}\cdot\fvec_{\bm{i}}\,h^d,
  \qquad
  (\Uvec,\Fvec)_{\mathcal U,h}
  =\sum_m\Uvec_m\cdot\Fvec_m\,w_m.
\]
Substituting the matched spreading formula and exchanging finite sums gives
\[
  \boxed{
  (\uvec,\mathcal S_h\Fvec)_h
  =(\mathcal J_h\uvec,\Fvec)_{\mathcal U,h}.}
\]

\begin{punchbox}
\centering
Thus \(\mathcal J_h=\mathcal S_h^{*}\) with respect to these inner products: the power delivered to the fluid equals the power removed from the structure.
\end{punchbox}

\medskip
This identity depends on the paired operators, their weights, and evaluation locations. It does \emph{not} require the sum-of-squares condition.
\end{frame}

% ----- Kernel catalog transition -----
\begin{frame}{Discrete kernels used in practice}
All of the following start from
\[
  \delta_h(\xvec)
  =\frac{1}{h^d}\prod_{\alpha=1}^{d}
  \phi\!\left(\frac{x_\alpha}{h}\right),
\]
but their one-dimensional functions \(\phi\) make different compromises.

\begin{itemize}\itemsep7pt
  \item Three-point kernels use a smaller stencil. \soft{(The Roma--Peskin--Berger kernel drops Peskin's even--odd condition and is meant for staggered MAC grids.)}
  \item Four-point kernels can combine exact low moments with greater smoothness.
  \item Simple analytic kernels may satisfy only a subset of Peskin's conditions.
  \item Wider kernels can target additional moments, smoother derivatives, or improved grid invariance.
\end{itemize}

\begin{punchbox}
Compare kernels property by property; labels such as ``higher order'' or ``smoother'' do not determine performance for every immersed problem.
\end{punchbox}
\end{frame}

% ----- 17. Kernel comparison -----
\begin{frame}{Three classical one-dimensional kernels}
\small
Write \(M_0\) and \(M_1\) for the exact zeroth- and first-moment conditions.
\begin{center}
\begin{tabular}{@{}lcccccc@{}}
\toprule
Kernel & Support & \(M_0\) & \(M_1\) & Even--odd & \(C\) in \(\sum\phi^2=C\) & Class \\
\midrule
Roma--Peskin--Berger 3-point & 3 & yes & yes & \warn{no} & \(C=1/2\) & \(C^1\) \\
Peskin 4-point & 4 & yes & yes & yes & \(C=3/8\) & \(C^1\) \\
Cosine 4-point & 4 & yes & \warn{no} & yes & \(C=3/8\) & \(C^1\) \\
\bottomrule
\end{tabular}
\end{center}
{\scriptsize Even--odd: \(\sum_{j\,\mathrm{even}}\phi(r-j)=\sum_{j\,\mathrm{odd}}\phi(r-j)=\tfrac12\) (deck 1-2). The three-point kernel drops it and is intended for staggered (MAC) grids (Peskin 2002, Sec.~6; Roma et al.\ 1999).\par}

\medskip
For the cosine kernel,
\[
  \phi_{\cos}(r)=
  \begin{cases}
  \dfrac14\left(1+\cos\dfrac{\pi r}{2}\right),& |r|\le2,\\[4pt]
  0,& |r|>2.
  \end{cases}
\]

\begin{punchbox}
The cosine kernel \hi{does satisfy} the constant sum-of-squares identity, but it fails the exact first-moment condition for generic subgrid offsets. These are different tests.
\end{punchbox}
\end{frame}

% ----- Wider and higher-order kernels -----
\begin{frame}{Higher-order and wider-support kernels}
Examples include six-point IB kernels (e.g.\ Bao, Kaye, and Peskin 2016) and B-spline-based families.

\begin{columns}[T,onlytextwidth]
  \begin{column}{0.48\textwidth}
  \begin{block}{Possible benefits}
  \begin{itemize}
    \item additional exact moments;
    \item smoother force and velocity transfer;
    \item reduced subgrid-position sensitivity;
    \item improved behavior for selected derivatives.
  \end{itemize}
  \end{block}
  \end{column}
  \begin{column}{0.48\textwidth}
  \begin{block}{Possible costs}
  \begin{itemize}
    \item a larger interaction stencil;
    \item more work and parallel communication;
    \item broader regularization of stress jumps;
    \item properties that may conflict, such as compactness and smoothness.
  \end{itemize}
  \end{block}
  \end{column}
\end{columns}

\begin{punchbox}
A wider kernel is not automatically more accurate for a solution with a singular interface force.
\end{punchbox}
\end{frame}

% ----- 18. Choosing a kernel -----
\begin{frame}{Kernel choice is an application-dependent tradeoff}
There is no universally best kernel. A defensible choice reports which properties the application needs.

\begin{columns}[T,onlytextwidth]
  \begin{column}{0.48\textwidth}
  \begin{block}{Numerical questions}
    \begin{itemize}
      \item Which moments are reproduced exactly?
      \item How smooth must interpolated motion and force be?
      \item Is nonnegativity important?
      \item What stencil width is affordable?
    \end{itemize}
  \end{block}
  \end{column}
  \begin{column}{0.48\textwidth}
  \begin{block}{Resolution questions}
    \begin{itemize}
      \item Are marker or quadrature points dense enough relative to \(h\)?
      \item Does the structure remain watertight as it deforms?
      \item How sensitive are results to subgrid translation and orientation?
    \end{itemize}
  \end{block}
  \end{column}
\end{columns}

\begin{punchbox}
Compare kernels with a problem-relevant convergence and conservation study; a wider or higher-moment kernel is not automatically more accurate near a singular interface.
\end{punchbox}
\end{frame}

% ----- 19. Refinement -----
\begin{frame}{What happens under refinement?}
For a normalized family with width proportional to \(h\),
\[
  \delta_h\longrightarrow\delta
  \qquad\text{in the sense of distributions as }h\to0.
\]

\begin{block}{That statement is about the kernel---not the whole algorithm}
Convergence of a complete IB discretization additionally requires:
\begin{itemize}
  \item consistent Eulerian and Lagrangian spatial discretizations,
  \item stable time integration and compatible refinement of \(\Delta t\),
  \item adequate marker or interaction-quadrature density,
  \item sufficient solution regularity for the claimed norm and order.
\end{itemize}
\end{block}

\begin{punchbox}
Regularization spreads an interface force over \(O(h)\), so stress jumps are smoothed at finite resolution. Leakage or residual slip need not disappear if the Lagrangian interaction points are under-resolved.
\end{punchbox}
\end{frame}

% ----- 20. Takeaways -----
\begin{frame}{Takeaways}
\begin{enumerate}
  \item The delta interaction is a \hi{push-forward from material measure to spatial measure}; it is not a hidden change of variables.
  \item A volumetric strong form contains both \(\operatorname{Div}\Pten\) in \(\mathcal U\) and the transmission traction \(-\Pten\Nvec\) on \(\partial\mathcal U\).
  \item Unified and partitioned weak forms organize those same two physical contributions differently.
  \item Zeroth and first moments control total force, torque, and low-order interpolation; constant sum of squares addresses approximate translation invariance.
  \item Matched spreading and interpolation should satisfy \(\mathcal J_h=\mathcal S_h^*\), which gives the discrete power identity.
\end{enumerate}

\begin{punchbox}
\centering
Always state the \hi{measure}, \hi{kernel properties}, and \hi{discrete inner products} together.
\end{punchbox}
\end{frame}

% ----- 21. References -----
\begin{frame}{Primary references}
\small
\begin{thebibliography}{99}
\bibitem{Peskin2002}
C.~S. Peskin,
``The immersed boundary method,''
\emph{Acta Numerica} 11 (2002), 479--517.
\href{https://doi.org/10.1017/S0962492902000077}{doi:10.1017/S0962492902000077}.

\bibitem{GriffithLuo2017}
B.~E. Griffith and X.~Y. Luo,
``Hybrid finite difference/finite element immersed boundary method,''
\emph{Int. J. Numer. Meth. Biomed. Eng.} 33 (2017), e2888.
\href{https://doi.org/10.1002/cnm.2888}{doi:10.1002/cnm.2888}.

\bibitem{GriffithPatankar2020}
B.~E. Griffith and N.~A. Patankar,
``Immersed methods for fluid--structure interaction,''
\emph{Annu. Rev. Fluid Mech.} 52 (2020), 421--448.
\href{https://doi.org/10.1146/annurev-fluid-010719-060228}{doi:10.1146/annurev-fluid-010719-060228}.

\bibitem{Roma1999}
A.~M. Roma, C.~S. Peskin, and M.~J. Berger,
``An adaptive version of the immersed boundary method,''
\emph{J. Comput. Phys.} 153 (1999), 509--534.
\href{https://doi.org/10.1006/jcph.1999.6293}{doi:10.1006/jcph.1999.6293}.
\end{thebibliography}
\vspace{2pt}
\soft{\scriptsize Adapted from: spreading/interpolation and kernel postulates after Peskin (2002), Secs.~5--6; the strong and weak elastic-force forms (\(\operatorname{Div}\Pten\) and \(-\Pten\Nvec\)) after Griffith and Luo (2017). Additional references at the end of the deck.}
\end{frame}

% =====================================================================
% Optional appendix: detailed derivations and kernel formulas
% =====================================================================
\section{Appendix}

\begin{frame}{Appendix: testing force on a spatial region}
Let \(\Gamma\subset\Omega\). A formal indicator-function calculation suggests
\[
  \int_{\Gamma}\fvec(\xvec,t)\,d\xvec
  =
  \int_{\{\Xvec\in\mathcal U:\,\chivec(\Xvec,t)\in\Gamma\}}
  \Fvec(\Xvec,t)\,d\Xvec.
\]

\begin{block}{The distributional caveat}
The rigorous starting point is the smooth-test-field identity. Approximate the indicator of \(\Gamma\) by smooth functions and pass to the limit. The region identity is unambiguous when the pushed-forward force measure assigns zero mass to \(\partial\Gamma\).
\end{block}

\begin{punchbox}
With a regularized \(\delta_h\), classical Fubini applies under ordinary integrability assumptions, but a marker within \(O(h)\) of \(\partial\Gamma\) contributes only partially to the region.
\end{punchbox}
\end{frame}

\begin{frame}{Appendix: Peskin's four-point kernel}
\small
\[
\phi_4(r)=
\begin{cases}
\dfrac18\left(3-2|r|+\sqrt{1+4|r|-4r^2}\right),
&0\le |r|<1,\\[6pt]
\dfrac18\left(5-2|r|-\sqrt{-7+12|r|-4r^2}\right),
&1\le |r|<2,\\[6pt]
0,&|r|\ge2.
\end{cases}
\]

\begin{columns}[T,onlytextwidth]
  \begin{column}{0.48\textwidth}
  \begin{itemize}
    \item support across four grid points
    \item zeroth and first moments exact
  \end{itemize}
  \end{column}
  \begin{column}{0.48\textwidth}
  \begin{itemize}
    \item \(\sum_j\phi_4(r-j)^2=3/8\)
    \item nonnegative and \(C^1\)
  \end{itemize}
  \end{column}
\end{columns}
\end{frame}

\begin{frame}{Appendix: Roma--Peskin--Berger three-point kernel}
\small
\[
\phi_3(r)=
\begin{cases}
\dfrac13\left(1+\sqrt{1-3r^2}\right),
&|r|<\dfrac12,\\[7pt]
\dfrac16\left(5-3|r|-\sqrt{-2+6|r|-3r^2}\right),
&\dfrac12\le |r|<\dfrac32,\\[7pt]
0,&|r|\ge\dfrac32.
\end{cases}
\]

\begin{punchbox}
This \(C^1\), nonnegative kernel satisfies the zeroth and first moment conditions and \(\sum_j\phi_3(r-j)^2=1/2\), with support across three grid points for generic offsets.
\end{punchbox}
\soft{\small It does \emph{not} satisfy the even--odd condition (at \(r=0\) the even sum is \(\phi_3(0)=2/3\)); Peskin (2002, Sec.~6) obtains it by dropping that condition and recommends it only with schemes, such as the staggered MAC scheme, that lack the \(D^0\) null-space problem.}
\end{frame}

\begin{frame}{Appendix: checking the cosine kernel}
\setlength{\abovedisplayskip}{5pt}\setlength{\belowdisplayskip}{5pt}%
For
\[
  \phi_{\cos}(r)=
  \begin{cases}
  \dfrac14\left(1+\cos\dfrac{\pi r}{2}\right),&|r|\le2,\\[4pt]
  0,&|r|>2,
  \end{cases}
\]
direct summation gives, for all real \(r\),
\[
  \sum_j\phi_{\cos}(r-j)=1,
  \qquad
  \sum_j\phi_{\cos}(r-j)^2=\frac38.
\]
However,
\[
  \sum_j(r-j)\phi_{\cos}(r-j)\ne0
\]
for a generic offset \(r\) (computed explicitly on the next slide).

\begin{punchbox}
The example shows why zeroth moment, first moment, and constant sum of squares must be checked separately.
\end{punchbox}
\end{frame}

\begin{frame}{Appendix: the cosine kernel's first moment, explicitly}
For \(0\le r<1\) the four contributing offsets are \(s=r+1,\ r,\ r-1,\ r-2\). With \(\theta=\pi r/2\), the cosines are \(-\sin\theta,\ \cos\theta,\ \sin\theta,\ -\cos\theta\), so
\[
  \sum_s s=4r-2,\qquad \sum_s s\cos\frac{\pi s}{2}=2\cos\theta-2\sin\theta,
\]
and therefore
\[
  \boxed{M_1(r)=\sum_j(r-j)\phi_{\cos}(r-j)=r-\tfrac12+\tfrac12\Bigl(\cos\tfrac{\pi r}{2}-\sin\tfrac{\pi r}{2}\Bigr).}
\]
\begin{itemize}\itemsep4pt
  \item \(M_1\) vanishes only at \(r=0,\ \tfrac12,\ 1\) and is odd about \(r=\tfrac12\).
  \item Its extremes are \(|M_1|\approx0.021\) (in units of \(h\)) at \(r\approx0.21\) and \(0.79\), where \(\sin\theta+\cos\theta=4/\pi\).
  \item The same bookkeeping with the squares gives \(\tfrac1{16}(4+2)=\tfrac38\), confirming the sum-of-squares value.
\end{itemize}
\end{frame}

\begin{frame}[plain]
\centering
\vfill
{\heading{Questions?}\par}

\medskip
{\color{DarkGray}\rule{0.34\textwidth}{0.8pt}}

\medskip
{\small Integral transforms and regularized delta functions}
\vfill
\end{frame}

% ----- Final references (appendix, kernel catalog, figures) -----
\begin{frame}{References (continued)}
\small
\begin{itemize}\itemsep5pt
  \item C.~S. Peskin, ``Numerical analysis of blood flow in the heart,'' \emph{J. Comput. Phys.} 25 (1977), 220--252.
  \href{https://doi.org/10.1016/0021-9991(77)90100-0}{doi:10.1016/0021-9991(77)90100-0}.
  \item C.~S. Peskin (2002), \emph{Acta Numerica} 11, Sec.~6: four-point kernel (eq.~6.27), cosine approximation (eq.~6.28), and the three-point variant; full entry under ``Primary references.''
  \item A.~M. Roma, C.~S. Peskin, and M.~J. Berger (1999): three-point kernel used with the staggered MAC scheme; full entry under ``Primary references.''
  \item F.~H. Harlow and J.~E. Welch, ``Numerical calculation of time-dependent viscous incompressible flow of fluid with free surface,'' \emph{Phys. Fluids} 8 (1965), 2182--2189 (MAC grid).
  \href{https://doi.org/10.1063/1.1761178}{doi:10.1063/1.1761178}.
  \item Y. Bao, J. Kaye, and C.~S. Peskin, ``A Gaussian-like immersed-boundary kernel with three continuous derivatives and improved translational invariance,'' \emph{J. Comput. Phys.} 316 (2016), 139--144 (six-point kernel).
  \item Figures: spreading and interpolation schematics from deck 1-2 (instructor's course notes).
\end{itemize}
\end{frame}

\end{document}
